%% notes-tex/interaction-fitness-figure-gate/sections/1-decision.tex
%% SOURCE: every number is a macro of tables/gate_numbers.tex, written by
%% experiments/030-solid-growth-multi/scripts/figure2_gate_numbers.py from the result files
%% of experiments 030 (lower_order_baselines_030.py, job 3902), 025
%% (additive_baselines_025.py), 010, 006, 005 and 007. The earlier documents named here are
%% notes-tex directories.

\section{The decision\secstatus{todo}}
\label{sec:decision}

Figure 2 of the manuscript is to carry fitness and gene--gene interactions, with gene
essentiality moved to its own figure. Every arm here is scored on the pinned held-out
triples of the Kuzmin 2018 and 2020 trigenic screens on the 030 build: arm R, the 010
random split over records (\gnRTrain{} train, \gnRVal{} validation, \gnRTest{} test
triples), and arm Q, the query-pair-disjoint split (\gnQTrain{}, \gnQVal{}, \gnQTest{}),
in which no held-out triple's query pair appears in any training record. Nothing here is
approved. The sheets are mockups of what the result files support, and the symbol after
each panel letter says only whether the values behind it are on file.

\pfstep{Where the full record is} \file{025-additive-baselines} holds the gene-identity
ladder on both arms, the arms and their missing control, the prespecified decision table
and the first arm Q test score of a transformer checkpoint. \file{025-s3-closure} holds the
closure: the published trigenic score reproduced from the stored lower-order records, and
what that identity says about training on the closure pool. \file{025-mmli-campaign} and
\file{030-scaling-study} hold the transformer campaigns on the 025 and 030 builds. The
lower-order ladder of this document is the script
\texttt{lower\_order\_baselines\_030.py} and its note.

\pfstep{The question the figure answers, and the one it must not dodge} The joint model
of experiment 030 trains on every single, every closure double and every training triple
of the S3 pool, with a fitness head beside the interaction head. For a held-out triple the
published identity
\begin{equation}
  \tau_{ijk} = f_{ijk} - f_{ij}\,f_k - \varepsilon_{ik} - \varepsilon_{jk}
  \label{eq:identity}
\end{equation}
has every term but $f_{ijk}$ in the training pool: the query strain's double fitness
$f_{ij}$, the array single $f_k$ and the two digenic control scores. So the model is handed
the triple's null expectation, and the question a reviewer will ask is whether its score is
that arithmetic. The gene-identity ladder of 010 and 025 cannot answer it, because a
one-hot ridge over the three genes never sees a measurement. The lower-order ladder can:
every rung reads only those terms, under the trainer's own label policy, and on arm Q the
terms of any held-out query pair are masked (\gmaskQFijq{} rows lose $f_{ij}$ and
\gmaskQEpsij{} lose $\varepsilon_{ij}$), exactly the doubles the strict closure rule
removed from the model's training pool.

\pfstep{Apples to apples, stated precisely} The information set is the same on both
sides: the lower-order rungs and the transformer see the identical training records, and
neither sees the held-out triple's own fitness or interaction, since the split is by
record and a 030 record holds all of a genotype's entries. What differs is the access
mode, and the sheets name it as a regime. A lower-order rung is a lookup plus a function:
the measured terms are inputs at scoring time, which is transductive over the closure.
The transformer as trained reads gene identity, graphs and sequence embeddings at scoring
time and must have stored the lower-order records in its weights. So the ladder is not a
claim that boosted trees are the better model of a cell. It is the measurement of how
much of a held-out trigenic score is already in the lower-order records, and the model's
distance from that number is the thing the figure has to show. A transformer that reads
the same terms as inputs is the arm that makes the two regimes directly comparable, and
it is the first entry of the registry (\ref{sec:registry}).

\pfstep{CGT is a fitness predictor, and that is the other fair reading} The joint model
predicts a held-out triple's fitness as well as its score, and on arm R its fitness head
reads \gfitRCgtValWindow{} against \gloRFOneridgelowerorderfitnesstest{} for the ridge
over the terms. Equation~\ref{eq:identity} then gives two rungs that use CGT as what it
is. C1 takes the model's predicted $f_{ijk}$ and subtracts the measured null,
$f_{ij} f_k + \varepsilon_{ik} + \varepsilon_{jk}$, the same terms the lower-order rungs
read, so the model supplies only the one term the closure lacks; it is the comparison in
which the model and the trees receive the same information in the same form. C2 derives
$\tau$ from the model's own predicted fitness of the triple, the doubles and the singles,
with no measured term, which tests whether its fitness predictions are consistent
enough for the residual. Neither is on file: both need the fitness head's per-triple
predictions, which E3 writes beside the interaction head's. If C1 reads above the trees,
the model's fitness head is doing the work the ladder measures and the figure says so;
if it does not, the fitness head is no more precise on $f_{ijk}$ than a ridge over the
terms, which Figure~\ref{fig:sheet}f already suggests on arm Q.

\pfstep{Three statements are on file} First, the free arithmetic is large. The
zero-parameter null that predicts $\tau$ as minus the two digenic control scores reads
\gloRNOnemultquerydoubletest{} on arm R and \gloQNOnemultquerydoubletest{} on arm Q (test
Pearson), boosted trees over the thirteen terms and four products
\gloRLTwohgblowerordertest{} and \gloQLTwohgblowerordertest{}, against the gene-identity
ridge at \gloRBOneadditivegenetest{} and \gloQBOneadditivegenetest{}; the paired gain of
the trees over the ridge is \gpairRBOneadditivegeneLTwohgblowerorder{}
[\gpairRBOneadditivegeneLTwohgblowerorderLo{}, \gpairRBOneadditivegeneLTwohgblowerorderHi{}]
on R and \gpairQBOneadditivegeneLTwohgblowerorder{}
[\gpairQBOneadditivegeneLTwohgblowerorderLo{}, \gpairQBOneadditivegeneLTwohgblowerorderHi{}]
on Q over 2,000 resamples of the test triples. Second, the transformer does not collect
it. The composite-embedding model with the fitness head and the dataset token reads
\gcgtRValWindow{} on arm R (validation, mean of epochs \gcgtRFirstEpoch{} to
\gcgtRLastEpoch{}, seed \gcgtRSeed{}) and \gcgtQValWindow{} on arm Q (epochs
\gcgtQFirstEpoch{} to \gcgtQLastEpoch{}), under the zero-parameter null on both arms, and
on Q within a screen it reads close to the screen-mean floor. Third, the fitness side is
arithmetic. A ridge over the same terms predicts a held-out triple's fitness at
\gloRFOneridgelowerorderfitnesstest{} on R and \gloQFOneridgelowerorderfitnesstest{} on
Q, and the model's fitness head reads \gfitRCgtValWindow{} and \gfitQCgtValWindow{}
(validation) on the same triples, so the fitness panel of the figure is a panel about
the closure, not about the model.

\pfstep{Every pooled number carries a floor} The two screens differ in mean $\tau$, and
predicting only the training mean of the triple's screen reads
\gloRNZeroscreenmeantest{} on R and \gloQNZeroscreenmeantest{} on Q. The dataset token
learns that offset in the first epoch. Every comparison on the sheets is therefore drawn
with the floor visible, and the performance table (\ref{sec:performance}) carries the
within-screen Pearson, both series centered by screen, beside the pooled one.

\pfstep{The claims the main-text figure would make} First, that the trigenic score is a
residual over lower-order measurements, and the ontology stores those measurements beside
it, which is what makes the closure recomputable and the free arithmetic measurable: on
file. Second, that a model reading gene identity and sequence recovers part of a held-out
triple's score on a random split and less on a query-pair-disjoint split: on file for one
seed, validation only. Third, that the model beats what the lower-order records already
say: not on file, and on the current reading false on both arms. Fourth, that the figure
can show the comparators of the earlier draft (DCell, DANGO, FBA) on the same data: not on
file; their results are on other builds (\ref{sec:models}). The figure enters the paper
only when the third claim is settled either way by the registry's first two entries, and
with the regime named in the panel.
